\documentclass[../../main.tex]{subfiles}

\begin{document}

\chapter{Polynomials}


\section{Foundations}

Before we get into the world of polynomials in olympiad algebra,
let's discuss some notations that we will be using throughout the
notes and key theorems as a review. Please feel free to skip over
to the next section if you are already familiar with the topics
below.

\begin{definition}{}{}
A \textbf{monomial} is a single term composed with variables raised
to powers of whole numbers and a coefficient. \textbf{Polynomial},
denoted as $P(x)$, is a sum of finitely many monomials.
\end{definition}

Some examples of monomial includes $x, 3a^2, 9, (5-2i)x^3y^2$ while
$x^{-2}, a + b, \sin{x}$ are not monomials. Polynomials are
typically written with the following sum notation.
\[
P(x)
= \sum_{i=0}^n a_i x^i
= a_n x^n + a_{n-1} x^{n-1} + \cdots + a_1 x + a_0
\]
With this in mind, we can define some notations that occasionally
appear.

\begin{definition}{}{}
Below are common set notations for polynomials where $n$ is a
non-negative integer.
\vspace{-0.8\baselineskip}
\begin{align*}
    \mathbb{Z}[x]
    &\coloneqq \left\{
        \sum_{i=0}^n a_i x^i \mid a_i \in \mathbb{Z}
    \right\} \\
    \mathbb{Q}[x]
    &\coloneqq \left\{
        \sum_{i=0}^n a_i x^i \mid a_i \in \mathbb{Q}
    \right\} \\
    \mathbb{R}[x]
    &\coloneqq \left\{
        \sum_{i=0}^n a_i x^i \mid a_i \in \mathbb{R}
    \right\} \\
    \mathbb{C}[x]
    &\coloneqq \left\{
        \sum_{i=0}^n a_i x^i \mid a_i \in \mathbb{C}
    \right\}
\end{align*}
\end{definition}

Continuing with the notations, let's discuss some key theorems
that cannot be omitted when discussing polynomials.

\begin{theorem}
{Fundamental Theorem of Algebra}
{Fundamental Theorem of Algebra}
Any nonconstant polynomial $P(x) \in \mathbb{C}[x]$ has at least one
complex root \cite{MATH01-3}.
\end{theorem}

There are multiple proofs for this beautiful theorem, but I think
they are out of scope for the notes. Trust me, the theorem is true,
and maybe I will include a rigorous proof of the statement in other
parts of LibreNotebook. Now this theorem can lead us to another very
beautiful proven assumption that we can use when solving problems.

\begin{corollary}{}{}
All polynomials $P(x) \in \mathbb{C}[x]$ with $\deg(P) = n > 1$
can be factored into $n$ linear factors with multiplicity.
\end{corollary}
\begin{proof}
Notice that by FTA, $P_n(x)$ can be represented as $P_n(x) =
a (x - z_1) P_{n-1}(x)$ where $P_k(x)$ is a polynomial with degree
$k$. Continuing recursively until $n = 1$, $P_n(x)$ can be written as
the following.
\begin{align*}
    P_n(x)
    &= a (x - z_1) (x - z_2) \cdots (x - z_{n-1}) P_1 (x) \\
    &= a (x - z_1) (x - z_2) \cdots (x - z_{n-1}) (x - z_n)
\end{align*}
Therefore, any polynomial with degree $n > 1$ can be factored into
exactly $n$ linear factors with multiplicity.
\end{proof}

The corollary also implies that there exists $n$ roots counting
multiplicity. Another key property to remember is remainder
theorem and factor theorem. They are pretty self-evident, so I
will omit the proof here, but below is the simple statement.

\begin{theorem}{}{}
For any polynomial $P$, if the polynomial is divided by $x - a$,
then the remainder is $P(a)$. Moreover if $P(a) = 0$, then $(x - a)$
is a factor of $P$.
\end{theorem}

The theorem above is basically stating that polynomials can be written
as the following.
\[
P(x) = (x - a) Q(x) + P(a)
\]
This is pretty self evident as we get $P(a) = P(a)$ if we substitute
$x = a$ and $Q(x)$ is the quotient polynomial.

Before we finish our review, here is the final theorem.

\begin{theorem}{Rational Root Theorem}{Rational Root Theorem}
Every rational root $\frac{p}{q}$ in lowest terms of a polynomial
$P(x) = \sum_{i=0}^n a_i x^i$ with $a_0, a_n \neq 0$ satisfies the
following conditions: $p \mid a_0$ and $q \mid a_n$.
\end{theorem}
\begin{proof}
First and foremost, notice that if $\frac{p}{q}$ is a solution to
$P(x) = 0$, then the following equations hold.
\begin{align*}
    P \left( \frac{p}{q} \right)
    = \sum_{i=0}^n a_i \left( \frac{p}{q} \right)^i
    &= 0 \\
    a_0 q^n + a_1 p q^{n-1} + a_2 p^2 q^{n-2} + \cdots + a_n p^n
    &= 0
\end{align*}
Therefore, the following equations can be derived.
\begin{align*}
    q \left( a_0 q^{n-1} + a_1 p q^{n-2} + a_2 p^2 q^{n-3} +
        \cdots + a_{n-1} p^{n-1}
    \right)
    &= -a_n p^n \\
    p \left( a_1 q^{n-1} + a_2 p q^{n-2} + a_3 p^2 q^{n-3} +
        \cdots + a_n p^{n - 1}
    \right)
    &= -a_0 q^n
\end{align*}
Because $p$ and $q$ are coprime, the first equation reveals that
$q \mid a_n$ and the second equation shows that $p \mid a_0$.
\end{proof}

This theorem sometimes comes handy if we had to blindly guess
a solution to further factorize as it narrows down our choices.
Now with the reminders, let's discuss the topic that we all may
have heard of: Roots of Unity.

\end{document}


